% sections/evaluation.tex — Evaluation Framework
\section{Evaluation Framework}
\label{sec:evaluation}

\system{} includes a built-in evaluation suite that validates the statistical,
structural, and domain-specific quality of generated data.  This section
describes the evaluation dimensions and the auto-tuning mechanism.

% ─────────────────────────────────────────────────────────────────────
\subsection{Statistical Validation}
\label{sec:eval_statistical}

\paragraph{Benford's Law Tests.}
First-digit and first-two-digit distributions are tested against theoretical
Benford probabilities using:
\begin{itemize}
  \item \textbf{Mean Absolute Deviation (MAD):}
    $\MAD = \frac{1}{K}\sum_{k=1}^{K}|\hat{p}_k - p_k|$, with conformity
    thresholds from~\citet{nigrini2012benfords}: close ($<0.006$), acceptable
    ($<0.012$), marginally acceptable ($<0.015$).
  \item \textbf{Chi-squared test:}
    $\chi^2 = n\sum_{k} (\hat{p}_k - p_k)^2 / p_k$.
  \item \textbf{Kolmogorov-Smirnov test} for the continuous approximation.
\end{itemize}

\paragraph{Distribution Fitting.}
Amount distributions are tested for log-normal fit using the
Anderson-Darling statistic at configurable significance levels (default
$\alpha = 0.05$).

\paragraph{Correlation Verification.}
The Pearson correlation matrix of generated fields is compared to the
configured target matrix.  Per-pair deviations exceeding the significance
threshold are flagged.

\paragraph{Temporal Pattern Analysis.}
Intraday activity profiles, day-of-week seasonality, and month-end
concentration ratios are compared against configured targets.  CUSUM and
EWMA control charts detect unexpected drift or regime changes.

% ─────────────────────────────────────────────────────────────────────
\subsection{Coherence Validators}
\label{sec:eval_coherence}

\system{} implements 15 domain-specific coherence validators, each verifying
cross-table consistency within a business domain:

\begin{table}[H]
\centering
\caption{Coherence validators and the invariants they enforce.}
\label{tab:validators}
\small
\begin{tabular}{p{3.2cm}p{9cm}}
\toprule
\textbf{Validator} & \textbf{Key Invariants} \\
\midrule
Balance Sheet &
  Assets $=$ Liabilities $+$ Equity; period-over-period consistency \\
Document Chain &
  PO$\to$GR$\to$Invoice$\to$Payment referential integrity; timestamp
  ordering \\
Subledger Reconciliation &
  AR/AP/FA/Inventory subledger totals $=$ GL control accounts \\
Intercompany Elimination &
  IC receivables $=$ IC payables; elimination entries net to zero \\
HR / Payroll &
  Gross $=$ Net $+$ Deductions; tax withholding rates within bounds \\
Manufacturing &
  BOM costing $\leq$ production order cost; scrap accounting \\
Sales Quotes &
  Quote$\to$order conversion consistency; discount limit adherence \\
Sourcing &
  Bid evaluation completeness; scorecard metric ranges \\
Bank Reconciliation &
  Book balance $+$ outstanding items $=$ bank balance \\
Financial Reporting &
  Statement interconnection (IS $\to$ BS via retained earnings;
  BS $\to$ CF via working capital changes) \\
Standards (IFRS/GAAP) &
  ASC 606 obligation satisfaction; ASC 842 ROU asset amortization;
  Fair value level 1+2+3 $=$ total \\
Network Analysis &
  Vendor concentration $\leq$ configured max; customer segment
  revenue shares within tolerance \\
Tax &
  Tax rate application; withholding accuracy \\
Treasury &
  Cash position flow consistency; covenant compliance \\
Project Accounting &
  Budget vs.\ actual variance within thresholds \\
\bottomrule
\end{tabular}
\end{table}

Each validator returns a structured result with pass/fail status, violation
count, violation details, and suggested remediation.

% ─────────────────────────────────────────────────────────────────────
\subsection{ML-Readiness Assessment}
\label{sec:eval_ml}

For datasets intended for machine learning, \system{} evaluates:

\paragraph{Feature Quality.}
Distribution matching between original fingerprint (if available) and
synthetic data using Wasserstein distance.  Categorical cardinality
verification.  Missing-data pattern analysis (MCAR/MAR/MNAR).

\paragraph{Label Quality.}
Anomaly label distribution across categories.  Label-feature correlation
analysis to verify that labels are predictable from features.  Ground-truth
completeness metrics.

\paragraph{Split Metrics.}
When temporal train/validation/test splits are applied, the evaluator checks
that: (a)~class distributions are balanced across splits, (b)~no future
data leaks into training, (c)~stratification preserves minority class
representation.

\paragraph{GNN Readiness.}
For graph exports, additional metrics assess:
\begin{itemize}
  \item Node feature dimensionality and scaling properties.
  \item Edge density and sparsity level.
  \item Graph connectivity (number of connected components, average degree).
  \item Heterogeneity: number of node types, edge types, and cross-type
    edges.
\end{itemize}

\paragraph{Scheme Detectability.}
For AML typology data, the evaluator computes per-typology detection scores
assessing whether the generated patterns have sufficient salience for ML
models while remaining realistic.

% ─────────────────────────────────────────────────────────────────────
\subsection{Extended Evaluators}
\label{sec:eval_extended}

Four additional evaluators were introduced in v0.11.0:

\paragraph{Multi-Period Coherence.}
Validates that opening balances equal prior closing balances, that document
IDs are sequential across periods, and that entity references remain
consistent throughout the generation horizon.

\paragraph{Fraud Pack Effectiveness.}
Measures detection rate at configurable thresholds, performs false positive
analysis, and computes coverage metrics per fraud scenario pack
(\Cref{sec:anom_fraud_packs}).

\paragraph{OCEL Enrichment Quality.}
Assesses state transition coverage, correlation event linking accuracy, and
resource utilization distribution across workload pools in OCEL~2.0 event
logs.

\paragraph{Causal Intervention Magnitude.}
Compares KPI deltas against expected magnitudes from the causal DAG, verifies
propagation paths, and validates that preservation constraints are maintained
under intervention.

Additionally, counterfactual evaluation is supported via the
\texttt{DiffEngine}, which computes record-level diffs between baseline and
counterfactual outputs (see Section~\ref{sec:counterfactual}).

% ─────────────────────────────────────────────────────────────────────
\subsection{AutoTuner}
\label{sec:eval_autotuner}

The AutoTuner closes the loop between evaluation and generation.  Given
evaluation results with identified gaps, it:

\begin{enumerate}
  \item Analyzes each gap to determine root cause (insufficient samples,
    distribution mismatch, weak correlations, etc.).
  \item Generates a prioritized list of \texttt{TuningOpportunity} records,
    each specifying a configuration patch.
  \item Ranks opportunities as Critical, Important, or Nice-to-have.
\end{enumerate}

Supported tuning actions include:
\texttt{IncreaseTransactionVolume},
\texttt{AdjustDistributionParameters},
\texttt{EnhanceAnomalyInjection},
\texttt{ImproveTemporalPatterns},
\texttt{BalanceFeatureDistributions},
\texttt{EnhanceGraphConnectivity}, and
\texttt{AdjustCorrelationMatrix}.

The output is a YAML patch that can be merged with the original
configuration for a subsequent generation run, enabling iterative
refinement toward quality targets.
